% Compile with: latexmk -pdf paper_v2.tex
\documentclass[10pt]{article}
\usepackage[margin=1in]{geometry}
\usepackage{times}
\usepackage{amsmath,amssymb,amsthm}
\usepackage{booktabs}
\usepackage{graphicx}
\usepackage{microtype}
\usepackage[authoryear,round]{natbib}
\usepackage{xcolor}
\usepackage{hyperref}

\newtheorem{proposition}{Proposition}
\newtheorem{definition}{Definition}
\newcommand{\todo}[1]{\textcolor{red}{\textit{[#1]}}}
\newcommand{\ess}{\mathrm{ESS}}
\newcommand{\E}{\mathbb{E}}
\newcommand{\R}{\mathbb{R}}

\title{How Many In-Context Examples Is a Task Description Worth?}
\author{Bruce Quan Nguyen \\ \todo{Affiliation}}
\date{Draft of 1 October 2026}

\begin{document}
\maketitle

\begin{abstract}
Prompts usually pair a task description with in-context examples, yet theories of in-context learning (ICL) model only the examples. We ask how many examples a description is worth. Reading ICL as Bayesian inference over a latent task, we treat the description and the examples as two kinds of data about the task, and define the description's effective sample size (ESS) as the number of examples that lower a Bayes-optimal predictor's one-step log-loss regret by as much as the description does. In in-context linear regression, where this predictor is exact, the ESS is governed by two properties of the description: its precision, the fraction $r$ of prior variance it leaves, and its reliability, the probability $p$ that it is correct. A reliable description is worth about $(d-1)(1-r) + 1/(r a_0)$ examples in $d$ dimensions at prior signal-to-noise ratio $a_0$. An unreliable description keeps at least a fraction $1-p$ of the regret of predicting from the prior alone, so its ESS saturates: at $p = 0.9$ in sixteen dimensions, a hundredfold gain in precision adds ten examples, against a hundred for a reliable description. Examples lift this cap by verifying the description, and the ESS then approaches $p$ times that of a reliable one; the regime that matters is a precise but imperfect description with few examples. In one training run per setting, small meta-trained Transformers are Bayes-optimal within their output class: they reproduce the ESS of a reliable description, and with an unreliable description alone, where the Bayes-optimal predictive is a two-component mixture, they match the best single Gaussian instead. Their trust is inherited from training: a network trained only on correct descriptions trusts an unreliable one fully, and a precise one then hurts it. Our claims concern Bayes-optimal predictors and small meta-trained models, not large language models.
\end{abstract}

\section{Introduction}

\begin{figure}[t]
\centering
\includegraphics[width=\linewidth]{figs/overview.pdf}
\caption{\textbf{A task description is worth a number of in-context examples, set by its precision and capped by its reliability.}
(a) We read a prompt as Bayesian inference over a latent task $w$ under the pretraining prior: the description is a direct noisy observation of $w$ and the examples are indirect ones. The description's effective sample size (ESS) is the number of examples $k$ that lowers a Bayes-optimal predictor's one-step regret by as much as the description does.
(b) Three descriptions and their ESS. The precision $r$ sets how narrowly the prior concentrates around the stated $m$; the reliability $p$ is the probability that the description is correct, and otherwise $w$ follows the base prior (dashed; its share $1-p$ shaded pink). Prior shapes are sketches; ESS values are from Table~\ref{tab:reliability} in Appendix~\ref{app:tables} ($d = 16$, $a_0 = 10$).}
\label{fig:overview}
\end{figure}

In-context learning (ICL) lets a model adapt to a new task from its prompt alone \citep{brown2020}. In practice the prompt usually pairs a natural-language description of the task with a few demonstrations, but the evidence on what the description contributes is mixed. A good description can be worth hundreds of labelled data points \citep{lescao2021, honda2025}, a misleading one can perform about as well as an informative one \citep{webson2022}, and enough examples override an explicitly stated bias \citep{gupta2025}. Formal theories of ICL do not settle the question, because they model the prompt as a set of demonstrations \citep{garg2022, akyurek2023, xie2022}, and the theoretical models that add a descriptor \citep{huangge2025, lin2025, tong2026} do not measure its value on the same scale as the examples. This leaves a basic question open: how many in-context examples is a task description worth?

We answer it within the Bayesian reading of ICL \citep{xie2022}. The pretraining distribution is a prior over a latent task, and everything in the prompt is data about that task: the examples observe it indirectly through their inputs, and the description observes it directly, if imperfectly. Conditioning on the description first gives the prior under which the examples are read, which is the sense in which ``the description is the prior''. Bayesian statistics already measures the information in a prior in units of observations, its \emph{effective sample size} (ESS) \citep{morita2008}, and finds that it depends on how concentrated the prior is and on whether it conflicts with the data \citep{evans2006}. For a task description these become two properties: its \textbf{precision}, how far it narrows the prior, and its \textbf{reliability}, the probability that it is correct.

We study the question where it can be answered exactly: noisy in-context linear regression with a Gaussian task prior \citep{garg2022, akyurek2023, raventos2023}. The setting buys two things. The Bayes-optimal predictor is available in closed form for every prompt, so a description's ESS is computed rather than estimated from a trained model. And the reliability of the descriptions a network is trained on is a knob we can turn, so the trust a meta-trained Transformer learns can be read off against two exact references, the predictor calibrated to the reliability and the one that trusts the description fully \citep{genewein2025}. A description is a noisy observation of the regression weights, and an unreliable description follows the robust mixture prior that clinical trials use to admit possibly conflicting historical data \citep{schmidli2014}.

\paragraph{Contributions.}
\begin{itemize}
    \item We define the ESS of a task description by matching the one-step log-loss regret of a Bayes-optimal predictor given the description to that given examples alone, with or without further examples already in the prompt. It is the prior sample size of \citet{reimherr2021} with predictive regret as the measure of uncertainty (\S\ref{sec:setting}).
    \item We show that a reliable description has $\ess \approx (d-1)(1-r) + 1/(r a_0)$: the share of the dimension it pins down at high signal-to-noise ratio, plus the conjugate-prior ESS of its own prior. We prove the expansion in the coarse and precise regimes, with error terms, and a bound that holds everywhere (\S\ref{sec:precision}).
    \item We show that an unreliable description's regret is at least $1-p$ times that of predicting from the prior alone, and at most that plus an identification cost of $H(p)$, so its ESS saturates however precise it is. Examples verify the description and lift the cap: the ESS with $n$ examples in hand approaches $p$ times that of a reliable description. A learner that trusts the description fully is instead hurt by precision (\S\ref{sec:reliability}).
    \item We meta-train small Transformers on the same process and place each between reference predictors. They are Bayes-optimal within their output class: they reproduce the ESS of a reliable description and, with an unreliable one alone, match the best single Gaussian rather than the Bayes-optimal mixture their head cannot express, merging with it as examples arrive. Their trust is inherited from training: trained only on reliable descriptions they behave like the fully trusting learner, and trained on unreliable ones they keep hedging when every description is correct (\S\ref{sec:networks}).
\end{itemize}
Together these results say that reliability is not a fixed discount on a description's value but a cost that examples pay down.

\section{Related Work}

\paragraph{Bayesian perspectives on ICL.}
ICL can be framed as implicit Bayesian inference over a latent task \citep{xie2022}. In linear regression, meta-trained Transformers implement Bayes-optimal ridge regression \citep{akyurek2023, raventos2023, panwar2024}. \citet{genewein2025} score meta-trained predictors by their expected excess log loss over the oracle that knows the task, and compare that regret with the Bayes-optimal predictor's; we adopt both the measure and the comparison. \citet{zhu2026} show that a Transformer given a prefix of related datasets performs amortised hierarchical Bayesian prediction and matches the oracle across prior families, the closest demonstration that a prior-carrying prefix is used Bayes-optimally. In contrast to all of these, we express the value of the conditioning information in examples, and we allow it to be wrong.

\paragraph{Task descriptions in ICL theory.}
\citet{huangge2025} add a descriptor token carrying the input mean; it helps their learner by centring the inputs but carries no information about the weights, so to a Bayes-optimal predictor its ESS is zero. \citet{lin2025} guide models with hypothesis classes, and \citet{tong2026} show that a prediction's dependence on an instruction, right or wrong, decays exponentially in the number of demonstrations. \citet{lin2024} analyse Gaussian-mixture task priors whose components are fixed by pretraining; ours has two components, one centred by the description. None of these works measures a description's value as an effective sample size.

\paragraph{Effective sample size.}
\citet{morita2008} define the ESS of a prior by matching its information to that of a baseline prior plus a sample. \citet{reimherr2021} make this prior sample size a function of the data in hand, allow it to be negative for a prior that conflicts with the likelihood, and leave prediction-based measures of uncertainty open; our ESS is their prior sample size with one-step predictive regret as that measure. \citet{neuenschwander2020} note that a mixture prior's ESS is not the weighted average of its components'. \citet{reznik2026} defines a borrowed ESS under predictive log loss for a power prior, with a divergence between historical and current data in the role our $1-p$ plays; it has no mixture and does not depend on the current data. Our model of reliability is the robust mixture prior of \citet{schmidli2014}, whose worked example puts weight $0.1$ on a vague component, our $p = 0.9$, with our base prior in the vague component's role. In that literature's terms a wrong description is a prior-data conflict \citep{evans2006}.

\section{Problem Formulation}\label{sec:setting}

\paragraph{Tasks and examples.}
A task is a weight vector $w \in \R^d$. An in-context example is an input $x \sim \mathcal{N}(0, I_d)$ with label $y = w^\top x + \epsilon$, $\epsilon \sim \mathcal{N}(0, \sigma_y^2)$, and the base prior over tasks is $w \sim \mathcal{N}(0, s_0^2 I_d)$. We set $\sigma_y = 1$ and write $a_0 = s_0^2/\sigma_y^2$ for the prior signal-to-noise ratio (SNR). We take $a_0 = 10$, above the values of $1$ to $4$ in the noisy settings of \citet{garg2022, raventos2023, panwar2024} and within the range swept by \citet{akyurek2023}, so that a precise description can be worth many examples.

\paragraph{Descriptions.}
A description is a vector $m \in \R^d$, generated with a latent indicator $Z \sim \mathrm{Bernoulli}(p)$ of whether it is correct:
\begin{equation}\label{eq:description}
m \sim \mathcal{N}\big(0, (1-r)s_0^2 I_d\big), \qquad
w \mid m, Z=1 \sim \mathcal{N}\big(m, r s_0^2 I_d\big), \qquad
w \mid m, Z=0 \sim \mathcal{N}\big(0, s_0^2 I_d\big).
\end{equation}
The marginal law of $w$ is the base prior in both cases, and the prior given the description is the robust mixture of \citet{schmidli2014},
\begin{equation}\label{eq:mixture}
p\,\mathcal{N}\big(m, r s_0^2 I_d\big) + (1-p)\,\mathcal{N}\big(0, s_0^2 I_d\big),
\end{equation}
with the base prior as its vague component. The description's \emph{precision} is measured by $r \in (0,1)$, the ratio of the variance around the description to the prior variance, so a smaller $r$ is a more precise description; the description locates each weight to within $\sqrt{r}$ of its typical size $s_0$, that is, $32\%$ at $r = 0.1$, $10\%$ at $r = 0.01$, and $3\%$ at $r = 0.001$. Its \emph{reliability} is $p$; we call a description \emph{reliable} if $p = 1$ and \emph{unreliable} otherwise.

Two features of this model matter below. First, a correct description is data about the task: equivalently, $m \mid w \sim \mathcal{N}((1-r)w,\; r(1-r)s_0^2 I_d)$, a direct noisy observation of $w$, so the ESS is an exchange rate between direct observations of $w$ and indirect ones through $x$. Second, a wrong description is unrelated to the task rather than deliberately misleading. A description that pointed away from $w$ by a fixed rule would be information about $w$ to a learner that knew the rule, and would harm only one that did not. Throughout, the description is a numerical observation, not text to be parsed.

\paragraph{Regret.}
A predictor sees a prompt $D$, which holds the description (if any), $n$ examples, and a query $x$, and outputs a predictive density $q(y \mid D)$. Its one-step regret is its expected excess log loss over the oracle that knows $w$,
\begin{equation}\label{eq:regret}
R = \E\big[-\log q(y \mid D) + \log \phi(y;\, w^\top x, \sigma_y^2)\big],
\end{equation}
where $\phi(\cdot\,; \mu, v)$ is the Gaussian density and the expectation is over tasks, descriptions, inputs, and labels: the per-step version of the regret of \citet{genewein2025}. We use log loss rather than the squared error of \citet{garg2022} because it is a proper scoring rule and the loss under which the mixture bound of \S\ref{sec:reliability} holds. We write $R^{\mathrm{desc}}(n)$ for the regret of the Bayes-optimal predictor given the description and $n$ examples, $R^{\mathrm{desc}}_{p=1}(n)$ for the same with a reliable description, and $R^{\mathrm{ex}}(n)$ for its regret given $n$ examples alone.

\paragraph{Effective sample size.}
\citet{reimherr2021} measure a prior $\pi$ against a baseline $\pi_b$ given data $I$ by the number $M$ of further observations that make the baseline as uncertain as the prior: $U_{\pi_b}(I + M) = U_\pi(I)$. We take the base prior as $\pi_b$, the description-conditioned prior~\eqref{eq:mixture} as $\pi$, the $n$ examples as $I$, and regret as $U$.

\begin{definition}[Effective sample size of a description]\label{def:ess}
Fix $p$ and $r$. The description's ESS with $n$ examples in hand is $\ess_n = n^* - n$, where $n^*$, interpolating $R^{\mathrm{ex}}$ linearly between integers, solves $R^{\mathrm{ex}}(n^*) = R^{\mathrm{desc}}(n)$. Its \emph{effective sample size} is $\ess = \ess_0$, its ESS alone.
\end{definition}

$\ess_n$ is the number of further examples an examples-only predictor needs to match the predictor that also has the description; it is negative when the description hurts. Because both sides are Bayes-optimal predictors with different conditioning sets, the ESS measures the information in the description rather than the architecture of any particular learner.

\paragraph{Two reference predictors.}
A learner need not know the reliability $p$. We write $q$ for the \emph{trust} it places in the description, the weight its posterior gives the component in which the description is correct, and two predictors recur. The \emph{calibrated} predictor sets $q = p$ and is the Bayes-optimal predictor above. The \emph{fully trusting} predictor sets $q = 1$: it takes every description as correct, as a learner that has never seen a wrong one would. Both are Bayes-optimal under their own belief about reliability. For any learner, its \emph{implied trust} is the $q$ whose predictor makes posterior-mean predictions closest to the learner's, in mean squared difference over prompts; it locates the learner between the two references.

\section{The Value of Precision}\label{sec:precision}

We first take a reliable description. Given $n$ examples $X_n \in \R^{n \times d}$, the posterior covariance of $w$ under a Gaussian prior with SNR $a$, that is, with variance $a\sigma_y^2$ per coordinate, is $\Sigma = \sigma_y^2 (I/a + X_n^\top X_n)^{-1}$, which depends on neither the labels nor the prior mean. The one-step regret is therefore
\begin{equation}\label{eq:gaussian-regret}
R = \tfrac12\, \E \log\big(1 + x^\top \Sigma x\big),
\end{equation}
half the expected log of the predictive variance left over, in units of the noise. A reliable description is the case $n = 0$, $a = r a_0$; examples alone have $a = a_0$.

\begin{proposition}[ESS of a reliable description]\label{prop:precision}
Let $p = 1$ and $c = 1/(r a_0) = \sigma_y^2/s_\ell^2$, where $s_\ell^2 = r s_0^2$. Then:
\begin{enumerate}
\item[(i)] \emph{(coarse descriptions)} if $c \le 1$, then as $rd \to \infty$, uniformly in $c \le 1$,
\[
\ess = (d-1)(1-r) + (1-r)\,c + O\big(1/(rd)\big);
\]
\item[(ii)] \emph{(precise descriptions)} if $c \ge d$, then as $c/d \to \infty$, uniformly in $a_0 \ge 1$,
\[
\ess = c + d + 1 - 1/a_0 + O(d/c);
\]
\item[(iii)] for every $d$, $a_0$, and $r$, $\ess \le c + d + 2$.
\end{enumerate}
\end{proposition}

\emph{Proof idea.} A description shrinks the prior variance by the factor $r$ in all $d$ directions. Examples remove variance differently. At high SNR each of the first $d$ examples removes the variance along the one direction it probes, so $n < d$ examples leave $d - n$ directions untouched; matching that leftover to the description's gives $d - n \approx rd$, the high-SNR term $(d-1)(1-r)$. Once $n \ge d$ every direction has been probed and the leftover decays like $d/(n - d - 1)$; matching it to the description's $d/c$ gives $n \approx c + d + 1$. The base prior has conjugate-prior ESS $1/a_0$ of its own, which both sides share; it is the $-1/a_0$ in (ii) and, through $(1-r)c = c - 1/a_0$, the correction in (i).

The constant $c = \sigma_y^2/s_\ell^2$ is the ESS of the description's own prior $\mathcal{N}(m, s_\ell^2 I_d)$ in the conjugate sense, noise variance over prior variance \citep{neuenschwander2020}, and \S\ref{sec:reliability} shows it is also the limit of $\ess_n$ as $n \to \infty$; we call it the conjugate-prior ESS. Both regimes say the same thing: the ESS of a reliable description is its high-SNR term plus its conjugate-prior ESS,
\begin{equation}\label{eq:additive}
\ess \approx (d-1)(1-r) + \frac{1}{r a_0},
\end{equation}
within about two examples: on the grid of Table~\ref{tab:precision} (Appendix~\ref{app:tables}) the rule is within $2.2$ examples of the exact ESS in every cell with $a_0 \ge 10$. The regimes meet at $c = 1$, where the description pins $w$ down as tightly as one label's noise does, and the ESS passes $d$ near $c = d$. Appendix~\ref{app:precision} gives the regret expansions behind (i) and (ii) and the constants in each regime.

In words, a description that removes nine tenths of the prior variance is worth nine tenths of the dimension plus one example, $14.9$ examples at $d = 16$ and $57.7$ at $d = 64$; one that leaves a thousandth of the variance is worth its conjugate-prior ESS of $100$ plus $d + 1$, or $117$ examples at $d = 16$. Figure~\ref{fig:ess}(a) traces the full curve at $d = 16$.

\section{The Bottleneck of Reliability}\label{sec:reliability}

With an unreliable description the Bayes-optimal posterior predictive is a mixture of two components, one that trusts the description and one that ignores it, weighted by the posterior probability that the description is correct. The following proposition bounds what this hedge costs, with or without examples to check the description against.

\begin{proposition}[Regret of an unreliable description]\label{prop:floor}
Fix $n \ge 0$, $p \in (0,1)$, and $r \in (0,1)$. Let $D_n$ be the prompt (the description, $n$ examples, and the query), $\pi_n = P(Z = 1 \mid D_n)$ the posterior probability that the description is correct, and $\pi_n(Z)$ the posterior weight on the true component, $\pi_n$ if $Z = 1$ and $1 - \pi_n$ otherwise. The predictor does no better than a twin who is told whether the description is correct, and no worse than that twin plus the cost of not knowing:
\begin{equation}\label{eq:sandwich}
p\,R^{\mathrm{desc}}_{p=1}(n) + (1-p)\,R^{\mathrm{ex}}(n)
\;\le\; R^{\mathrm{desc}}(n) \;\le\;
p\,R^{\mathrm{desc}}_{p=1}(n) + (1-p)\,R^{\mathrm{ex}}(n) + \E\big[-\log \pi_n(Z)\big].
\end{equation}
The identification term $\E[-\log \pi_n(Z)]$ is the conditional entropy $H(Z \mid D_n)$; it equals the binary entropy $H(p)$, in nats, at $n = 0$ and is non-increasing in $n$. In particular, alone the regret is at least $(1-p)\,R^{\mathrm{ex}}(0)$ for every $r$, so the ESS cannot exceed the $n$ at which $R^{\mathrm{ex}}(n) = (1-p)\,R^{\mathrm{ex}}(0)$; and as $r \to 0$ the regret alone lies between $(1-p)\,R^{\mathrm{ex}}(0)$ and $(1-p)\,R^{\mathrm{ex}}(0) + H(p)$.
\end{proposition}

\emph{Proof idea.} Compare the predictor with a twin that is also told $Z$. The twin's regret is the left side of~\eqref{eq:sandwich}, and the predictor cannot beat it, since extra information never hurts a Bayes-optimal predictor. Conversely, the predictor's posterior predictive puts weight $\pi_n(Z)$ on the twin's answer, so it loses at most $-\log \pi_n(Z)$; with no examples that is $\log(1/p)$ when the description is correct and $\log(1/(1-p))$ when it is wrong, which average to $H(p)$. More examples can only sharpen the posterior on $Z$ in expectation. The upper bound is the standard bound for a mixture forecaster under log loss, which pays at most the log of one over the weight it put on the right component \citep[Cor.~3.1 and \S9.2]{cesabianchi2006} \citep[\S V]{merhav1998}, and the lower bound is the Bayes-optimality of the informed twin; what is new is what the two bounds imply for the ESS. The proof is in Appendix~\ref{app:floor}.

\begin{figure}[t]
\centering
\includegraphics[width=\linewidth]{../../results/rq2-reliability/ess.pdf}
\caption{\textbf{Reliability caps precision only while nothing can check the description.} ESS of a description against its precision $r$ with $n$ examples in hand ($\ess_n$ of Definition~\ref{def:ess}), one panel per reliability; $d = 16$, $a_0 = 10$, log scales, one seed of 8{,}000 prompts. Dotted: the limit of the ESS alone as $r \to 0$, $135$ examples at $p = 0.99$ and $25$ at $p = 0.9$. Alone, the $p = 0.9$ description reaches $23$ examples at $r = 0.001$, close to that limit, while the reliable one keeps rising; with $100$ examples in hand the three panels nearly coincide ($109$, $107$, and $88$ at $r = 0.001$).}
\label{fig:ess}
\end{figure}

\paragraph{Alone, reliability caps the ESS.}
With no examples nothing in the prompt distinguishes the two components, since $m$ has the same law whether or not it is correct, so the predictor must hedge. Proposition~\ref{prop:floor} caps the ESS at the $n$ where $R^{\mathrm{ex}}(n) = (1-p)\,R^{\mathrm{ex}}(0)$: $40$ examples at $p = 0.9$ and $325$ at $p = 0.99$, for $d = 16$ and $a_0 = 10$. The actual limit is lower, because in the precise limit the identification term is all of $H(p)$: a description that pins $w$ down exactly or is unrelated to it leaves nothing to tell the two cases apart. As $r \to 0$ the ESS saturates at $25$ examples at $p = 0.9$ and $135$ at $p = 0.99$ (dotted in Figure~\ref{fig:ess}). The bounds are close to the computed regret: at $r = 0.001$ and $p = 0.9$ they are $0.32$ and $0.64$ nats and the regret is $0.57$, against $2.51$ with the prior alone and $0.07$ for a reliable description.

Table~\ref{tab:reliability} (Appendix~\ref{app:tables}) shows the saturation across precision. Refining $r$ from $0.1$ to $0.001$ raises the ESS of a reliable description from $14.9$ to $116$, but that of a description with $p = 0.9$ only from $13.1$ to $23.1$. A one-in-a-hundred chance of error nearly halves the ESS of the most precise description, $64$ against $116$, yet leaves coarse descriptions almost untouched: the $p = 0.99$ curve stays within a tenth of the reliable one down to $r = 0.01$. The same saturation appears at $d \in \{4, 64\}$ and at $a_0 = 100$.

\paragraph{Examples lift the cap.}
Examples can reveal whether the description is correct. In~\eqref{eq:sandwich} the lower bound $(1-p)\,R^{\mathrm{ex}}(n)$ falls with $n$ and the identification term cannot rise, so the cap rises as examples accumulate. Once the examples have identified the description and $R^{\mathrm{ex}}(n) \approx d/(2n)$, the reliable-description regret is about $d/(2(n + c))$, and setting $d/(2n^*)$ equal to the lower bound gives
\begin{equation}\label{eq:ess-n}
\ess_n \;\approx\; \frac{n\,p\,c}{n + (1-p)\,c} \;\xrightarrow[n \to \infty]{}\; p\,c = \frac{p}{r a_0},
\end{equation}
against $\ess_n \to c$ for a reliable description. Once verified, then, a precise unreliable description is worth $p$ times a reliable one. Identification requires the two components to be distinguishable: the identification term tends to the residual entropy of $Z$ given $w$ and $m$, which is negligible here because the components are far apart (their Kullback--Leibler divergence is of order $\tfrac{d}{2}\log(1/r)$) but stays positive for a coarse description, $0.17$ nats at $d = 5$, $r = 0.5$.

Before this limit the exact predictor passes through two regimes, shown in Figure~\ref{fig:ess}(c) for $r = 0.001$, $p = 0.9$, $d = 16$, and $c = 100$. The first example or two identify the description: the identification term falls from $H(p) = 0.33$ nats to $0.07$ after one example and $0.02$ after two, and the ESS rises from $23$ to $30$. The ESS then stays near $30$ while $n < d$, because the wrong-description component still carries the full examples-only regret, and rises once $n > d$ as that regret shrinks: $29$ at $n = 10$, $37$ at $n = 16$, and $88$ at $n = 100$, against $82$ from~\eqref{eq:ess-n} and a limit of $90$. Over the same range the reliable description's ESS drifts from $117$ to $108$, so at $n = 100$ the unreliable description is worth four fifths of the reliable one, where alone it was worth a fifth. A description with $c < d$ has little left to verify: at $r = 0.01$, where $c = 10$, the unreliable description gains one example after the first and then loses ESS like the reliable one.

\paragraph{A fully trusting learner is hurt by precision.}
The cap and its lifting assume a predictor calibrated to $p$. A predictor that takes the description as certainly correct, with trust $q = 1$ when the truth is $p = 0.9$, pays for the wrong tenth in full: its posterior predictive is as tight as the description, so when the description is wrong the answer falls far outside it. At $d = 16$ and $a_0 = 10$ the description alone costs this predictor $2.2$ nats at $r = 0.1$, against $2.51$ for no description, so a coarse description still helps ($6.5$ examples, against $13.2$ when calibrated). At $r = 0.01$ it costs $6.3$ nats and at $r = 0.001$ it costs $13.4$, worse than no description at all. Examples repair the damage only slowly, because the wrong tenth carries a prior as tight as the description: after $100$ examples the $r = 0.01$ description still leaves the predictor $58$ examples behind one without a description, and the $r = 0.001$ description leaves $4.0$ nats of regret against $0.09$. To a learner that cannot doubt the description, precision is a liability. This is the fully trusting predictor of \S\ref{sec:setting}; \S\ref{sec:networks} shows that a network trained only on correct descriptions behaves like it. Figure~\ref{fig:trust} draws the whole range of trust: the regret is flat for a wide band below $p$, rises gently above it, and jumps at $q = 1$; with ten examples in hand the band is flat everywhere except the jump, which examples do not remove.

\begin{figure}[t]
\centering
\includegraphics{../../results/rq2-reliability/trust.pdf}
\caption{\textbf{Under-trusting a description is cheap; trusting it fully is not.} Regret of the Bayes-optimal predictor against the trust $q$ it places in a description that is correct with probability $p = 0.9$, alone (solid) and with ten examples in hand (dashed), at $d = 16$, $a_0 = 10$; one seed of 8{,}000 prompts. At $r = 0.001$ alone, trusting $0.85$ costs $0.01$ nats more than the calibrated $0.566$ and trusting fully costs $12.5$; with ten examples every trust below $1$ gives $0.26$ and full trust still $10.8$.}
\label{fig:trust}
\end{figure}

\section{Meta-Trained Transformers}\label{sec:networks}

The results so far concern the Bayes-optimal predictor. A network meta-trained on prompts from the same process is not that predictor: it has an architecture, an output head, and a training distribution, and each can leave a mark. We ask two questions. Does a trained network value a description as the Bayes-optimal predictor does, and where it does not, what sets the difference? And what does a network learn about descriptions it was never shown wrong?

\paragraph{Setup.}
We meta-train Transformers on prompts from the process of \S\ref{sec:setting} and compare them with the Bayes-optimal predictor on the same prompts, following \citet{genewein2025}. Prompts use the prefix embedding of \citet{huangge2025}: an optional descriptor token carrying $m$, up to $20$ example tokens each holding an input beside its answer, and a final query token, with two indicator coordinates marking the token type. There is no positional encoding, so the examples are exchangeable. The architecture and optimiser follow \citet{garg2022}: $12$ layers, $8$ heads, width $256$, and Adam at a learning rate of $10^{-4}$, with fresh prompts at every step. The network outputs a mean and a log variance for the query's answer and is trained on the Gaussian log loss, so its regret is on the scale of~\eqref{eq:regret}. We set $d = 5$ and $a_0 = 10$ (the network sees $w \sim \mathcal{N}(0, I)$ and noise variance $0.1$, the same task in the units of \citet{garg2022}), draw the number of examples uniformly from $0$ to $20$, and omit the description in half the prompts. We train one network for each combination of a precise ($r = 0.01$) or coarse ($r = 0.1$) description with a reliable ($p = 1$) or unreliable ($p = 0.9$) one. The descriptor token supplies only $m$, so each network must learn $r$ and $p$ from its training distribution. Each trains for $40{,}000$ steps at batch $1024$ without a curriculum (\citet{garg2022} use batch $64$, $500{,}000$ steps, and a curriculum over dimension and prompt length), one training run per setting. In a pilot, the mean regret gap to the Bayes-optimal predictor fell from $0.07$ nats at $5{,}000$ steps to $0.03$ at $20{,}000$ and stayed near $0.02$ from $30{,}000$ to $40{,}000$.

\paragraph{Readouts and reference predictors.}
We evaluate each network on $20{,}000$ fresh prompts for every $n$ from $0$ to $20$, with and without a description, and on prompts of both reliabilities, so that each network meets descriptions it was trained on and descriptions it was not. As in Definition~\ref{def:ess}, the network's ESS is where its regret with the description alone meets its own regret curve for examples alone. We place each network between the reference predictors of \S\ref{sec:setting} on the same prompts: the \emph{calibrated} predictor, with trust equal to the test reliability, and the \emph{trained} predictor, with trust equal to the training reliability, which is the predictor each network is trained toward and becomes the fully trusting predictor when a network trained at $p = 1$ meets unreliable descriptions. A third reference belongs to the architecture: the \emph{best single Gaussian}, the predictive that matches the mean and variance of the Bayes-optimal mixture, which is all that a mean-and-variance output head can represent. We also read off each network's implied trust. The standard error of each regret is below $0.02$ nats, except for the precise $p = 1$ model tested at $p = 0.9$, where it reaches $0.10$ with a description at small $n$; the examples-only curve falls by about $0.13$ nats per example near $n = 0$, so these move an ESS by at most $0.15$ examples, and by $0.8$ in that last case.

\begin{table}[t]
\centering\small
\caption{Networks against the reference predictors ($d=5$, $a_0=10$, one training run per model, 20{,}000 evaluation prompts). The first four rows test each network on its training distribution, the last four on the other reliability. ``Trained'' is the Bayes-optimal predictor with trust equal to $p_{\mathrm{train}}$, ``calibrated'' the one with trust equal to $p_{\mathrm{test}}$ (shown where the two differ), and ``single Gaussian'' the best predictive a mean-and-variance head can express, shown where the Bayes-optimal predictive is a two-component mixture. The last two columns are the mean regret gap, network minus the trained predictor, over $n = 0, \dots, 20$, with and without a description.}
\label{tab:networks}
\input{../../results/rq3-meta-trained/networks_table.tex}
\end{table}

\begin{figure}[t]
\centering
\includegraphics[width=\linewidth]{../../results/rq3-meta-trained/regret.pdf}
\caption{\textbf{Networks are Bayes-optimal within their output class: they match the mixture where it is a single Gaussian and the best single Gaussian where it is not.} Regret against the number of examples for the network (markers; standard errors over 20{,}000 prompts are smaller than the markers) and the Bayes-optimal predictor (lines), one training run per panel. In each panel the lower curve is prompts with a description followed by $n$ examples and the upper curve is prompts with $n$ examples alone; the ESS is the $n$ at which the upper curve reaches the lower curve's value at $n = 0$. (a) $r=0.01$, $p=1$; (b) $r=0.1$, $p=1$; (c) $r=0.01$, $p=0.9$; (d) $r=0.1$, $p=0.9$. In (c) and (d) the Bayes-optimal predictive with a description is a two-component mixture; the dashed line is the best single Gaussian, and the network lies on it.}
\label{fig:networks}
\end{figure}

\paragraph{Networks are Bayes-optimal within their output class.}
On their training distribution the networks reproduce the Bayes-optimal ESS of a reliable description, $14.7$ against $15.4$ for the precise one and $5.24$ against $5.28$ for the coarse one, with regret $0.02$ to $0.03$ nats above the Bayes-optimal predictor's on average over $n$ (Table~\ref{tab:networks}, Figure~\ref{fig:networks}a, b). With an unreliable description alone they fall short: ESS $4.3$ against $6.8$ at $r = 0.01$ and $3.6$ against $4.2$ at $r = 0.1$, with regret $0.5$ and $0.2$ nats too high, and their implied trust reads $0.85$ rather than $0.9$. The shortfall is the output head, not the learning. With an unreliable description and no examples, the Bayes-optimal predictive is two Gaussians, one around the description's answer and one around the prior's, and a head that emits one mean and one variance cannot be both. The best it can do is the single Gaussian with the mixture's mean and variance, whose ESS is $4.31$ at $r = 0.01$ and $3.62$ at $r = 0.1$: the networks' $4.30$ and $3.57$. As examples identify the description the mixture collapses to one component and the three predictors merge (Figure~\ref{fig:networks}c, d): by $n = 5$ the best single Gaussian is within $0.01$ nats of the mixture and the network within its usual $0.05$, and the network's ESS with examples in hand rises from $4.3$ to $7.3$ after five and $8.1$ after ten (mixture $8.8$ and $9.8$), reproducing the lifting of \S\ref{sec:reliability}; at $r = 0.1$ it falls, from $3.6$ to $2.3$ and $1.3$, as the mixture's does. The implied trust of $0.85$ is what a unimodal predictive looks like through posterior means alone, not a belief the network holds about descriptions. Two controls support this reading: the shortfall is absent wherever the Bayes-optimal predictive is a single Gaussian (every reliable row, and every $n$ beyond a few examples), and it is not under-training, since a model trained at $r = 0.05$, $p = 0.9$ with the same recipe and continued to $80{,}000$ steps on fresh prompts kept its ESS ($3.9$ against $5.0$). With one training run per setting we state this as what the runs show, not as a general law.

\paragraph{Trust is inherited from training.}
The last four rows of Table~\ref{tab:networks} test each network on the reliability it was not trained on. A network trained only on correct descriptions behaves like the fully trusting learner of \S\ref{sec:reliability}: when one description in ten is wrong, the precise model's regret with the description alone is $2.87$ nats against $1.87$ without one, so a description hurts it, its ESS is zero where a calibrated predictor gets $6.9$, and its implied trust is $0.999$; the fully trusting Bayes-optimal predictor does worse still, $3.32$ nats. The coarse model is forgiven most of its trust, ESS $2.3$ against the calibrated $4.3$, because a wrong coarse description does little harm. The converse holds too: a network trained at $p = 0.9$ keeps hedging when every description is correct, with ESS $6.4$ against the calibrated $15.6$ at $r = 0.01$ and $4.6$ against $5.3$ at $r = 0.1$, and implied trust $0.85$ either way. Trust, then, is learned from the training distribution (in the pilot at $p = 1$, the implied trust alone rose from $0.9$ at $5{,}000$ steps to $0.99$ by $10{,}000$ and stayed there) and carried into deployment unchanged, and the exact regret as a function of trust (Figure~\ref{fig:trust}) says which error is the expensive one: at $d = 5$, $r = 0.01$, $p = 0.9$, trusting $0.85$ instead of $0.9$ costs $0.006$ nats and trusting fully costs $2.8$. A learner that has never seen a wrong description pays the full price of unreliability on precise descriptions; one that has seen them pays a small price for caution wherever they turn out to be reliable.

\section{Discussion}

For a Bayes-optimal predictor a task description has a definite value in in-context examples. A reliable description is worth the share of the dimension it pins down plus the conjugate-prior ESS $1/(r a_0)$, the examples that would match its variance. An unreliable description forces the predictor to hedge against its being wrong, so the value of a very precise description saturates. The hedge is forced only while nothing can check the description: a few examples identify it, and a precise unreliable description is then worth more than it was alone, approaching $p$ times a reliable one. Reliability is therefore not a fixed discount but a cost that examples pay down. This holds for a predictor calibrated to the description's reliability. To one that cannot doubt the description, as a network trained only on correct descriptions cannot, a precise description that is sometimes wrong is worse than none, and examples repair the damage slowly. Hedging also needs a predictive that can be two things at once: a network whose head emits one mean and one variance is Bayes-optimal within that class and pays for an unreliable description alone with the gap between the best single Gaussian and the mixture, which examples close. A language model's distribution over tokens can represent such a mixture; a regression head cannot.

\paragraph{Limitations.}
Our claims cover Bayes-optimal predictors and small meta-trained Transformers in a linear-Gaussian setting chosen for exact calculation; whether they carry over to large language models reading text is open. The network results rest on one training run per setting. The ESS scores one-step prediction: Definition~\ref{def:ess} counts the examples already in the prompt, but not a sequential session in which the predictor's own feedback accumulates. Over such horizons a reliable description's ESS tends to the number of examples whose expected information gain matches its $\tfrac{d}{2}\log(1/r)$ nats, below the one-step value ($7.8$ against $14.9$ at $d = 16$, $a_0 = 10$, $r = 0.1$); the horizon dependence of unreliable descriptions is left for future work.

\bibliographystyle{plainnat}
\bibliography{refs}

\clearpage
\appendix
\section{Exact values}\label{app:tables}

Tables~\ref{tab:precision} and~\ref{tab:reliability} give the exact one-step ESS behind \S\ref{sec:precision} and \S\ref{sec:reliability} at the ends of the grids; Figure~\ref{fig:ess} shows the relationships between them. The two tables come from separate simulations, so the one quantity they share, the reliable description at $d = 16$, $r = 0.001$, differs in its last digit: $117$ against $116$.

\begin{table}[h]
\centering\small
\caption{Exact ESS of a reliable description ($p=1$) for $a_0=10$, three seeds (standard deviation across seeds below 2\% of each value), against the rule $(d-1)(1-r) + 1/(r a_0)$ of~\eqref{eq:additive}, which is within $2.2$ examples of every cell.}
\label{tab:precision}
\input{../../results/rq1-single-query-gap/precision_table.tex}
\end{table}

\begin{table}[h]
\centering\small
\caption{ESS of unreliable descriptions ($d=16$, $a_0=10$), three seeds; standard deviation across seeds at most 2\% of each value. Down a column precision rises tenfold per row; across a row the description is wrong never, one time in a hundred, and one time in ten.}
\label{tab:reliability}
\input{../../results/rq2-reliability/reliability_table.tex}
\end{table}

\section{Proof of Proposition~\ref{prop:precision}}\label{app:precision}

We prove Proposition~\ref{prop:precision} in the following fuller form, which adds the regret expansions behind (i) and (ii); its final sentence is (iii).

\medskip\noindent\textbf{Proposition~\ref{prop:precision} (full statement).} {\itshape Let $p = 1$, let $\psi$ be the digamma function, and let $c = 1/(r a_0) = \sigma_y^2 / s_\ell^2$: the ESS of the description's prior $\mathcal{N}(m, s_\ell^2 I_d)$ in the conjugate sense, noise variance over prior variance \citep{neuenschwander2020}, the conjugate-prior ESS.
\begin{enumerate}
\item[(i)] \emph{Coarse descriptions.} Suppose $c \le 1$. Then for every $n < d$,
$R^{\mathrm{ex}}(n) \ge \tfrac12\big[\log(2a_0) + \psi\big(\tfrac{d-n}{2}\big)\big]$, with equality in the limit $a_0 \to \infty$, and as $rd \to \infty$, uniformly in $c \le 1$,
\[
\ess = (d-1)(1-r) + (1-r)\,c + O\!\left(\tfrac{1}{rd}\right).
\]
\item[(ii)] \emph{Precise descriptions.} Suppose $c \ge d$. Then for every $n \ge d$,
$R^{\mathrm{ex}}(n) \le \Phi(n) := \tfrac12\big[\psi\big(\tfrac{n+1}{2}\big) - \psi\big(\tfrac{n-d+1}{2}\big)\big]$, with equality in the limit $a_0 \to \infty$, and as $c/d \to \infty$, uniformly in $a_0 \ge 1$,
\[
\ess = c + d + 1 - \tfrac{1}{a_0} + O\!\left(\tfrac{d}{c}\right).
\]
\end{enumerate}
For every $d$, $a_0$, and $r$, $\ess \le c + d + 2$.}
\medskip

The two regimes say the same thing. In both regimes the ESS of a reliable description is its high-SNR term plus its conjugate-prior ESS,
$\ess \approx (d-1)(1-r) + 1/(r a_0)$,
to within about two examples. In (i) the conjugate term enters as $(1-r)c$, which falls short of $c$ by $1/a_0 < 1$; in (ii) the constant $c + d + 1 - 1/a_0$ exceeds the rule by $2 - 1/a_0 + (d-1)r \in [1, 2)$, up to the stated error. The regimes meet at $c = 1$, where the description pins $w$ down as tightly as one label's noise does, and the ESS passes $d$ near $c = d$. Below that point each gain in precision removes directions; above it every direction has already been probed, and the ESS grows like $c$, the examples needed to match the description's variance, plus the $d + 1$ that $d$ noisy examples cannot supply. The bound (iii) holds throughout, including the crossover $1 < c < d$ where neither expansion applies. On the grid of Table~\ref{tab:precision} (Appendix~\ref{app:tables}) the rule~\eqref{eq:additive} is within $2.2$ examples of the exact ESS in every cell with $a_0 \ge 10$; at $d = 64$, $a_0 = 100$, $r = 0.1$ the exact ESS is $56.85$ against $56.79$ from (i), and at $d = 16$, $a_0 = 10$, $r = 0.001$ it is $116.8$ against $116.9$ from (ii).

\newtheorem{lemma}{Lemma}[section]

\paragraph{Standard facts.}
Throughout, $\varepsilon = 1/a_0$, $A = \|x\|^2 \sim \chi^2_d$ for the query $x$, $\Sigma = (\varepsilon I + X_n^\top X_n)^{-1}$, and $g(t) = \E \log(1 + At)$, which is increasing and concave with
\begin{equation}\label{eq:g-derivs}
d\,(1 - (d+2)t) \le g'(t) = \E \frac{A}{1 + At} \le d,
\qquad
-(d^2 + 2d) \le g''(t) \le 0 .
\end{equation}
We use (a) $\E \log \chi^2_k = \psi(k/2) + \log 2$, $\E[\chi_k^{-2j}] = \prod_{i=1}^{j} (k - 2i)^{-1}$ for $k > 2j$, and $\psi(z) = \log z - 1/(2z) + O(z^{-2})$; (b) for $W \sim W_q(p, I)$ with $p \ge q$ and $z \sim \mathcal{N}(0, I_q)$ independent of $W$, $z^\top W^{-1} z$ is distributed as $\chi^2_q / \chi^2_{p-q+1}$ with independent numerator and denominator \citep[Theorem~3.2.12]{muirhead1982}, and $\E \operatorname{tr} W^{-2} = q(p-1)/((p-q)(p-q-1)(p-q-3))$ for $p - q \ge 4$ \citep{vonrosen1988}; (c) $y - y^2/2 \le \log(1 + y) \le y$ for $y \ge 0$.

\paragraph{Regret.}
The posterior predictive is the true conditional law of $y$, a Gaussian with variance $1 + x^\top \Sigma x$, so its expected log loss is its entropy and $R^{\mathrm{ex}}(n) = \tfrac12 \E \log(1 + x^\top \Sigma x)$, as in~\eqref{eq:gaussian-regret}. It is strictly decreasing in $n$, since by the Sherman--Morrison formula each new example lowers $x^\top \Sigma x$ almost surely. For the description, $n = 0$ and $\Sigma = r a_0 I = I/c$, so by (a) and (c) with $y = c/A$,
\begin{equation}\label{eq:rdesc}
2R^{\mathrm{desc}} = g(1/c) = \log(2 r a_0) + \psi(d/2) + \delta^{\mathrm{desc}},
\qquad
\delta^{\mathrm{desc}} = \E \log(1 + c/A) = \frac{c}{d-2} + O\!\left(\frac{c^2}{d^2}\right) \quad (d \ge 5).
\end{equation}

\paragraph{Locating the root.}
Let $h(n) = 2R^{\mathrm{ex}}(n) - 2R^{\mathrm{desc}}$ on integers $n \ge 0$, with linear interpolant $\bar h$. Then $\ess$ is the unique root of $\bar h$, and $h(0) = g(a_0) - g(r a_0) > 0$. In (i) and (ii) we show, for some $n_0 \ge 0$, $\alpha > 0$, and $\theta \to 0$ along the stated limit,
\begin{equation}\label{eq:root}
h(n) = \alpha\,(n_0 - n) + O(\alpha\theta)
\qquad \text{uniformly over integers $n \ge 0$ with $|n - n_0| \le 3$.}
\end{equation}
Being a convex combination of two such values, $\bar h(t)$ obeys the same expansion at real $t \ge 0$ with $|t - n_0| \le 2$; so for small $\theta$ the root lies within $2$ of $n_0$, and $0 = \bar h(\ess) = \alpha(n_0 - \ess) + O(\alpha\theta)$ gives $\ess = n_0 + O(\theta)$.

\begin{lemma}[fewer than $d$ examples]\label{lem:few}
Let $n < d$ and $k = d - n$. Then $2R^{\mathrm{ex}}(n) = \log(2a_0) + \psi(k/2) + \delta^{\mathrm{ex}}$ with $\delta^{\mathrm{ex}} \ge 0$, $\delta^{\mathrm{ex}} \to 0$ as $a_0 \to \infty$, and, for $k \ge 5$,
\[
\delta^{\mathrm{ex}} = \frac{\varepsilon (d-1)}{(k-1)(k-2)} + O\!\left(\frac{\varepsilon^2 d^2}{k^4}\right).
\]
\end{lemma}

\begin{proof}
Write $X_n^\top X_n = \sum_{i \le n} \lambda_i u_i u_i^\top$ with $\lambda_i > 0$ almost surely, and let $Q$ be the squared norm of the projection of $x$ onto the $k$ directions orthogonal to the $u_i$. Then $\Sigma$ has eigenvalue $a_0$ on those directions and $1/(\varepsilon + \lambda_i)$ on $u_i$, so
\[
1 + x^\top \Sigma x = a_0 Q + T, \qquad T = 1 + V, \quad V = \sum_{i \le n} \frac{z_i^2}{\varepsilon + \lambda_i}, \quad z_i = u_i^\top x,
\]
where, given $X_n$, the $z_i$ are i.i.d.\ standard normal and independent of $Q \sim \chi^2_k$. Taking logarithms and expectations, by (a),
\[
2R^{\mathrm{ex}}(n) = \log(2a_0) + \psi(k/2) + \delta^{\mathrm{ex}}, \qquad \delta^{\mathrm{ex}} = \E \log\!\left(1 + \frac{\varepsilon T}{Q}\right) \ge 0 .
\]
As $a_0 \to \infty$, $\varepsilon T$ decreases to $0$ pathwise, so $\delta^{\mathrm{ex}} \to 0$ by dominated convergence (the integrand at $a_0 = 1$ is integrable).

For the expansion, let $F = \sum_i z_i^2 / \lambda_i$. In the eigenbasis of $W = X_n X_n^\top \sim W_n(d, I)$, which has the same eigenvalues, $F = z^\top W^{-1} z$ and $\sum_i z_i^2/\lambda_i^2 = z^\top W^{-2} z$ with $z \sim \mathcal{N}(0, I_n)$ independent of $W$. By (b), $F \sim \chi^2_n / \chi^2_{k+1}$, so $\E F = n/(k-1)$ and $\E (1+F)^2 = O(d^2/k^2)$, and $\E \sum_i z_i^2/\lambda_i^2 = n(d-1)/(k(k-1)(k-3))$. Since $1/\lambda - \varepsilon/\lambda^2 \le 1/(\varepsilon + \lambda) \le 1/\lambda$, we have $F - \varepsilon \sum_i z_i^2/\lambda_i^2 \le V \le F$. Now (c) with $y = \varepsilon T / Q$, the independence of $Q$ from $T$, and (a) give
\[
\frac{\varepsilon\,(1 + \E V)}{k-2} - \frac{\varepsilon^2\, \E T^2}{2(k-2)(k-4)} \le \delta^{\mathrm{ex}} \le \frac{\varepsilon\,(1 + \E F)}{k-2} = \frac{\varepsilon (d-1)}{(k-1)(k-2)} ,
\]
and the two sides differ by $\varepsilon^2 \E \sum_i z_i^2/\lambda_i^2 /(k-2) + \varepsilon^2 \E T^2 / (2(k-2)(k-4)) = O(\varepsilon^2 d^2 / k^4)$, as $\E T^2 \le \E (1 + F)^2$.
\end{proof}

\begin{lemma}[at least $d$ examples]\label{lem:many}
Let $n \ge d$ and $m = n - d + 1$. Then $2\Phi(n) = \E \log(1 + A/Q)$ with $Q \sim \chi^2_m$ independent of $A$, and $R^{\mathrm{ex}}(n) \le \Phi(n)$, with equality as $a_0 \to \infty$. If moreover $a_0 \ge 1$ and $m \ge \max(d, 7)$, then
\[
2R^{\mathrm{ex}}(n) = 2\Phi(n) - \frac{\varepsilon d}{m^2} + O\!\left(\frac{\varepsilon d^2}{m^3}\right).
\]
\end{lemma}

\begin{proof}
Rotate coordinates so that $x = \|x\| e_1$; the rows of $X_n$ remain i.i.d.\ $\mathcal{N}(0, I_d)$, independent of $A$. Then $x^\top \Sigma x = A\,\Sigma_{11} = A/S$ with $S$ the Schur complement of the first coordinate in $\Sigma^{-1}$, which the Woodbury identity writes, with $X_1$ the first column of $X_n$ and $X_{-1}$ the other $d-1$ columns, as
\[
S = \varepsilon + X_1^\top \big(I_n + a_0 X_{-1} X_{-1}^\top\big)^{-1} X_1 .
\]
Given $X_{-1}$, the middle matrix has eigenvalue $1$ on the $m$ directions orthogonal to the columns of $X_{-1}$ and $1/(1 + a_0 \mu_j)$ on the other $d - 1$, where $\mu_j$ are the eigenvalues of $M = X_{-1}^\top X_{-1} \sim W_{d-1}(n, I)$. Since $X_1 \sim \mathcal{N}(0, I_n)$ is independent of $X_{-1}$, its coordinates in that eigenbasis are i.i.d.\ standard normal, and
\[
S = Q + \Delta, \qquad Q \sim \chi^2_m, \qquad
\Delta = \varepsilon + \sum_{j < d} \frac{z_j^2}{1 + a_0 \mu_j} \in \big[\varepsilon,\, \varepsilon(1 + F)\big], \qquad F = \sum_{j < d} \frac{z_j^2}{\mu_j},
\]
with $A$, $Q$, and $(z, M)$ independent; by (b), $F \sim \chi^2_{d-1} / \chi^2_{m+1}$, so $\E F = (d-1)/(m-1)$ and, when $m \ge \max(d, 7)$, $\E (1+F)^2 \le 5$.

Since $A + Q \sim \chi^2_{n+1}$, (a) gives $\E \log(1 + A/Q) = \psi\big(\tfrac{n+1}{2}\big) - \psi\big(\tfrac{m}{2}\big) = 2\Phi(n)$. As $S \ge Q$, $R^{\mathrm{ex}}(n) \le \Phi(n)$, with equality as $a_0 \to \infty$ by monotone convergence, since $\Delta$ decreases to $0$.

For the expansion, $2R^{\mathrm{ex}}(n) = 2\Phi(n) - \E D$ with
\[
D = \log\Big(1 + \frac{A}{Q}\Big) - \log\Big(1 + \frac{A}{S}\Big) = \log(1 + y_1) - \log(1 + y_2), \quad y_1 = \frac{\Delta}{Q} \ge y_2 = \frac{\Delta}{Q + A}, \quad y_1 - y_2 = \frac{\Delta A}{Q(Q+A)} .
\]
By concavity of $\log(1+y)$ and $1/(1+y_1) \ge 1 - y_1$, $(y_1 - y_2)(1 - \Delta/Q) \le D \le y_1 - y_2$. Here $\Delta$ is independent of $(A, Q)$; $\E \frac{A}{Q(Q+A)} = \E\frac1Q - \E\frac{1}{Q+A} = \frac{d}{(m-2)(m+d-2)}$ by (a), as $Q + A \sim \chi^2_{m+d}$; $\frac{A}{Q^2(Q+A)} \le \frac{A}{Q^3}$; and $\varepsilon \le \E\Delta \le \varepsilon(1 + \E F) = \varepsilon\frac{m+d-2}{m-1}$, $\E \Delta^2 \le 5\varepsilon^2$. Hence
\[
\frac{\varepsilon d}{(m-2)(m+d-2)} - \frac{5\,\varepsilon^2 d}{(m-2)(m-4)(m-6)} \le \E D \le \frac{\varepsilon d}{(m-1)(m-2)} ,
\]
and both sides are $\varepsilon d / m^2 + O(\varepsilon d^2 / m^3)$, using $\varepsilon \le 1$.
\end{proof}

\begin{proof}[Proof of Proposition~\ref{prop:precision}]
\emph{(iii).} Conditionally on $A$, $\log(1 + Av)$ is concave in $v$, so Jensen's inequality with $\E[1/Q] = 1/(m-2)$ gives $2\Phi(n) \le g(1/(m-2))$ for $m \ge 3$. If $m - 2 \ge c$, that is $n \ge c + d + 1$, then $g(1/(m-2)) \le g(1/c) = 2R^{\mathrm{desc}}$, so $R^{\mathrm{ex}}(n) \le \Phi(n) \le R^{\mathrm{desc}}$ by Lemma~\ref{lem:many}. The least such integer is $\lceil c + d + 1 \rceil < c + d + 2$, and $R^{\mathrm{ex}}$ is decreasing, so $\ess \le c + d + 2$.

\emph{(i).} The lower bound and the limit are Lemma~\ref{lem:few}. For the expansion let $c \le 1$, $n_0 = (d-1)(1-r) + (1-r)c$, and $n \ge 0$ an integer with $|n - n_0| \le 3$; write $k = d - n = rd + u$. Since $n_0 = d - rd - u_0$ with $u_0 = 1 - r + \varepsilon - c \in [-c, 1]$, we have $|u| \le 4$, so $k \ge 5$ once $rd \ge 9$. By Lemma~\ref{lem:few}, \eqref{eq:rdesc}, $\psi(z) = \log z - 1/(2z) + O(z^{-2})$, and $\varepsilon = rc$, with every $O$ uniform in $c \le 1$ and $|u| \le 4$,
\begin{align*}
h(n) &= \psi(k/2) - \psi(d/2) - \log r + \delta^{\mathrm{ex}} - \delta^{\mathrm{desc}} \\
&= \Big[\log\frac{k}{rd} - \frac1k + \frac1d\Big] + \frac{\varepsilon d}{r^2 d^2}\Big(1 + O\Big(\frac{1}{rd}\Big)\Big) - \frac{c}{d} + O\Big(\frac{1}{(rd)^2}\Big)
 = \frac{u - 1 + r + c - \varepsilon}{rd} + O\Big(\frac{1}{(rd)^2}\Big),
\end{align*}
since $\log(k/(rd)) = u/(rd) + O((rd)^{-2})$, $1/k = 1/(rd) + O((rd)^{-2})$, $\varepsilon d / (r^2 d^2) = c/(rd)$, and $c/d = \varepsilon/(rd)$. So $h(n) = \frac{1}{rd}\big[(n_0 - n) + O(1/(rd))\big]$, which is~\eqref{eq:root} with $\alpha = \theta = 1/(rd)$, and $\ess = (d-1)(1-r) + (1-r)c + O(1/(rd))$.

\emph{(ii).} The upper bound and the limit are Lemma~\ref{lem:many}. For the expansion let $n_0 = c + d + 1 - \varepsilon$ and $n$ an integer with $|n - n_0| \le 3$, so that $m - 2 = c + s$ with $s = n - n_0 - \varepsilon$, $|s| \le 4$; for $c \ge 9d$ every such $n$ has $m \ge \max(d, 7)$. By~\eqref{eq:rdesc} and Lemma~\ref{lem:many}, $2\Phi(n) - 2R^{\mathrm{desc}} = \E\big[g(1/Q) - g(1/c)\big]$, and a second-order Taylor expansion of $g$ about $1/c$ with~\eqref{eq:g-derivs}, $\E[1/Q] - 1/c = -s/(c(c+s))$, and $\E(1/Q - 1/c)^2 = \operatorname{Var}(1/Q) + s^2/(c(c+s))^2 = O(c^{-3})$ gives
\[
2\Phi(n) - 2R^{\mathrm{desc}} = -\frac{s}{c(c+s)}\,d\,\Big(1 + O\Big(\frac dc\Big)\Big) + O\Big(\frac{d^2}{c^3}\Big) = -\frac{ds}{c^2} + O\Big(\frac{d^2}{c^3}\Big).
\]
With Lemma~\ref{lem:many}, $\varepsilon \le 1$, and $\varepsilon d/m^2 = \varepsilon d/c^2 + O(d/c^3)$,
\[
h(n) = -\frac{d}{c^2}\Big[s + \varepsilon + O\Big(\frac{d}{c}\Big)\Big] = \frac{d}{c^2}\Big[(n_0 - n) + O\Big(\frac{d}{c}\Big)\Big],
\]
which is~\eqref{eq:root} with $\alpha = d/c^2$ and $\theta = d/c$, uniformly in $a_0 \ge 1$, so $\ess = c + d + 1 - 1/a_0 + O(d/c)$.
\end{proof}

\section{Proof of Proposition~\ref{prop:floor}}\label{app:floor}

\begin{proof}
Let $D_n = (x_{1:n}, y_{1:n}, m, x)$ be everything the predictor sees before predicting $y$, and $Z = 1$ if the description is correct, $Z = 0$ otherwise.

\emph{Step 1: the posterior predictive is a $\pi_n$-mixture.} The Bayes-optimal predictor's posterior predictive is $q(y) = \pi_n f_1(y) + (1 - \pi_n) f_0(y)$, where $\pi_n = P(Z = 1 \mid D_n)$ and $f_1$, $f_0$ are the predictives of the two components given $D_n$: $f_1$ is the reliable-description predictor's, and $f_0$ is the examples-only predictor's, because under $Z = 0$ the description is independent of $w$ and carries no information about $y$. At $n = 0$, $\pi_0 = p$: seeing $m$ says nothing about $Z$, because $m$ has the same law $\mathcal{N}(0, (1-r) a_0 I)$ either way, and $x$ is independent of $Z$, so there is nothing to update on. The components are then $f_1 = \mathcal{N}(m^\top x,\, 1 + r a_0\|x\|^2)$ and $f_0 = \mathcal{N}(0,\, 1 + a_0\|x\|^2)$.

\emph{Step 2: the twin.} Given $(D_n, Z)$, the true conditional density of $y$ is $f_Z$. So a twin who is told $Z$ predicts with $f_1$ when $Z = 1$, with regret $R^{\mathrm{desc}}_{p=1}(n)$, and with $f_0$ when $Z = 0$, with regret $R^{\mathrm{ex}}(n)$. Its regret is $p\,R^{\mathrm{desc}}_{p=1}(n) + (1-p)\,R^{\mathrm{ex}}(n)$.

\emph{Step 3: the predictor cannot beat the twin (lower bound).} Given $(D_n, Z)$, the expected log loss of any density $q$ exceeds that of the true density $f_Z$ by $\E[-\log q \mid D_n, Z] - \E[-\log f_Z \mid D_n, Z] = \mathrm{KL}(f_Z \,\|\, q) \ge 0$. Averaging over $Z$ and $D_n$ gives $R^{\mathrm{desc}}(n) \ge p\,R^{\mathrm{desc}}_{p=1}(n) + (1-p)\,R^{\mathrm{ex}}(n)$.

\emph{Step 4: the predictor loses at most the identification term (upper bound).} Pointwise, $q \ge \pi_n f_1$ and $q \ge (1 - \pi_n) f_0$, so $-\log q \le -\log f_Z - \log \pi_n(Z)$ with $\pi_n(1) = \pi_n$ and $\pi_n(0) = 1 - \pi_n$. Taking expectations gives the upper bound with the identification term $\E[-\log \pi_n(Z)]$. At $n = 0$ it is $p \log(1/p) + (1-p)\log(1/(1-p)) = H(p)$.

\emph{Step 5: the identification term is non-increasing.} $\E[-\log \pi_n(Z)] = H(Z \mid D_n)$, the conditional entropy of $Z$ given the data. The query $x$ is independent of $Z$ and of everything else in $D_n$, so $H(Z \mid D_n) = H(Z \mid x_{1:n}, y_{1:n}, m)$, and these conditioning sets are nested in $n$. Conditioning on more does not increase conditional entropy, so the term is non-increasing.

\emph{Step 6: the $n = 0$ case.} $R^{\mathrm{desc}}_{p=1}(0) = \tfrac12 \E \log(1 + r a_0 \|x\|^2)$ is nonnegative, which gives the lower bound $(1-p)\,R^{\mathrm{ex}}(0)$ for every $r$. The ESS is the $n$ at which $R^{\mathrm{ex}}(n) = R^{\mathrm{desc}}(0) \ge (1-p)\,R^{\mathrm{ex}}(0)$, and $R^{\mathrm{ex}}(n)$ is non-increasing in $n$ (a predictor given $n + 1$ examples can ignore one, and Bayes-optimal regret cannot rise with more information), so the ESS cannot exceed the $n$ at which $R^{\mathrm{ex}}(n) = (1-p)\,R^{\mathrm{ex}}(0)$. Finally $R^{\mathrm{desc}}_{p=1}(0)$ decreases to $0$ as $r \to 0$ by monotone convergence, which gives the two limits.
\end{proof}

\end{document}
